\scanpage{4}

\noindent\hfill Lodève le 27.\unsure{8.74|8/74|8/4}\\
Cher Deligne,\\
% margin note written sideways
\margin{Cher Breen, J'imagine que tu sais ce que le \unsure{maître|maitre} demande, mais n'ai pu \unsure{neamoins|neanmoins} lui repondre --\\
Qu'en \unsure{pens|penses}-tu. Bien à toi\\
P. \unsure{Deligne}\\}
Étant peut-être empêché \unsure{par ma|par une|par mon} \unsure{faute|faculté|famille} d'\unsure{assurer} un cours de 1\textsuperscript{er} cycle au 1\textsuperscript{er} trimestre, je vais peut\unsure{'-|-}être à la place faire un petit séminaire d'algèbre, et envisage de le faire sur les fourbis de Mme Sinh, éventuellement transposés dans le contexte des ``champs''. À ce \unsure{propos}, je tombe sur le truc suivant, qui pour l'instant reste heuristique. Si $M,N$ sont \unsure{deux} faisceaux abéliens sur un topos $X$, et $\tau_{\leq 2}\mathbb{R}\underline{\mathrm{Hom}}(M,N)=E(M,N)$ est le complexe ayant les invariants
\[
\left\{\begin{aligned}
\underline{H}^{\unsure{i|0}} &= \deleted{\unsure{\mathrm{Hom}}}\underline{\mathrm{Ext}}^{i}(M,N) && \text{pour}\quad 0\leq i\leq 2\\
\underline{H}^{\unsure{j|i}} &= 0 && \text{si}\quad i\notin[0,2]
\end{aligned}\right.
\]
il doit y avoir un triangle distingué canonique
\diagram{\begin{tikzcd}[ampersand replacement=\&,column sep=small,row sep=small]
\& \& {\underline{\mathrm{Hom}}(M,{}_{2}N)[-2]}\arrow[dl] \& \\
(T) \& {E(M,N)}\arrow[rr] \& \& {E'(M,N)}\arrow[ul]
\end{tikzcd}}

% archive sheet number: 61
